\section{RL Finetuning: de imitar trayectorias humanas a optimizar la calidad de la generación}

La clase ubica la falla principal de SFT en «from humans»: las soluciones humanas son costosas, tienen cobertura limitada y no necesariamente se encuentran en la región que el modelo genera con facilidad. La alternativa es que el modelo proponga sus propios candidatos, filtrarlos con un resultado final verificable y devolverle las trayectorias exitosas. Este ciclo une rejection sampling, self-training y, de forma más general, reinforcement learning finetuning.

\subsection{Primer paso: cambiar la fuente de datos por la distribución del propio modelo}

La slide 19 conserva primero el proceso de dos pasos de SFT y luego tacha «human-annotated»; las slides 20--21 muestran dos mejoras sucesivas: primero se maximizan las trayectorias correctas generadas por el modelo y después se reduce, al mismo tiempo, la probabilidad de las trayectorias incorrectas. Lo central no es que «los datos sintéticos siempre superan a los datos humanos», sino que los ejemplos de entrenamiento provienen de la política actual $\pi_\theta$ y, por lo tanto, se acercan más a los estados que el modelo realmente alcanza.

\lecturefigure{slide-19.jpg}{Corregir la falla de generalización de SFT: cuestionar primero las trayectorias humanas como fuente de datos}{19}

\lecturefigure{slide-20.jpg}{Ciclo de ejemplos positivos de RL finetuning: el modelo genera, se verifica y se vuelve a entrenar}{20}

\lecturefigure{slide-21.jpg}{Distinguir además entre trayectorias correctas e incorrectas para reforzar la preferencia relativa}{21}

Supongamos que el modelo genera $K$ trayectorias $y^{(k)}$ para el problema $x$ y que el verifier devuelve $v(x,y^{(k)})\in\{0,1\}$. La actualización más simple por rejection sampling conserva solo los ejemplos con $v=1$:
\[
\mathcal D_{\text{RS}}=\{(x,y^{(k)}):v(x,y^{(k)})=1\},
\qquad
\mathcal L_{\text{RS}}=-\mathbb E_{(x,y)\sim\mathcal D_{\text{RS}}}\log p_\theta(y\mid x).
\]
Aquí $K$ determina la amplitud de la exploración y $v$ determina qué trayectorias entran al conjunto de entrenamiento. Tras varias rondas, la distribución del modelo se desplaza hacia las regiones donde antes generó trayectorias exitosas; pero si el verifier se equivoca, las trayectorias incorrectas también se amplifican una y otra vez.

\begin{lstlisting}[language=Python,caption={Generación de datos por rejection sampling con análisis de la respuesta y verifier}]
def build_round(policy, problems, verifier, parse_answer, n=16):
    accepted, rejected = [], []
    for problem in problems:
        for sample in policy.sample(problem, n=n):
            answer = parse_answer(sample.text)
            if answer is None:
                rejected.append((problem, sample.text, "parse_error"))
            elif verifier(problem, answer):
                accepted.append((problem, sample.text))
            else:
                rejected.append((problem, sample.text, "wrong_answer"))
    return accepted, rejected
\end{lstlisting}

\teachervoice{00:20:00--00:24:50, el docente recuerda que también le sorprendió la primera vez que oyó que «las respuestas generadas por una máquina pueden ser más adecuadas para entrenar que los datos humanos». Experiencia práctica: la ventaja no proviene de que «la máquina sea más inteligente», sino de que los ejemplos están alineados con la distribución del modelo actual y de que es posible volver a muestrear después de cada mejora, lo que forma un ciclo iterativo cerrado.}

\subsection{Optimizar directamente lo que realmente queremos}

Las slides 22--24 condensan algoritmos complejos en una sola frase: primero se define la generation quality y luego se aplica gradient ascent. Para una política $\pi_\theta(y\mid x)$ y una recompensa $R(x,y)$, el objetivo puede escribirse como
\[
J(\theta)=\mathbb E_{x\sim\mathcal X,\,y\sim\pi_\theta(\cdot\mid x)}[R(x,y)],
\qquad
\nabla_\theta J
=\mathbb E[R(x,y)\nabla_\theta\log\pi_\theta(y\mid x)].
\]
Aquí $R$ puede ser la corrección de una respuesta matemática, la tasa de pruebas de código aprobadas, la calidad de una traducción o una puntuación de preferencia aprendida. La fórmula subraya que la recompensa determina la dirección de la optimización; policy gradient, las actualizaciones group-relative o el SFT con filtrado son solo métodos distintos para estimar e implementar esa dirección.

\lecturefigure{slide-22.jpg}{Por qué usar trayectorias generadas por el modelo: optimizar directamente la conducta objetivo}{22}

\lecturefigure{slide-23.jpg}{Primero se define la generation quality; el resto del trabajo consiste en calcular el gradiente}{23}

\lecturefigure{slide-24.jpg}{Abstracción de RL finetuning: ascenso por gradiente sobre la calidad esperada}{24}

\begin{importantbox}{La función objetivo va antes que el nombre del algoritmo}
Decir «usamos RL» no indica, por sí mismo, qué optimiza el sistema. Primero hay que especificar la reward, el análisis de la respuesta, la distribución de muestreo, el costo por longitud, el tratamiento de las salidas inválidas y las restricciones de KL o de estabilidad; solo después se discute PPO, GRPO, REINFORCE o rejection sampling. Cuanto mejor se optimiza un objetivo equivocado, más peligroso se vuelve el sistema.
\end{importantbox}

En el ejemplo de calidad, el material de la clase menciona una métrica de traducción automática, que normalmente debe escribirse BLEU. Las métricas son solo aproximaciones: el exact match en matemáticas puede ignorar expresiones equivalentes, las pruebas unitarias de código pueden pasar por alto casos límite ocultos y BLEU puede favorecer la coincidencia superficial de n-gramas. La brecha entre el objetivo real y la reward calculable es la puerta de entrada al reward hacking.

\subsection{El verifier es el cuello de botella actual de la escalabilidad automática}

La slide 25 cita «Verification, the key to AI» y ofrece un juicio explícito: un verifier confiable suele importar más que el algoritmo de RL específico. Sea $z\in\{0,1\}$ la corrección real y $v$ la salida del verifier. Si la false-positive rate
\[
\alpha=P(v=1\mid z=0)
\]
no es suficientemente baja, cuanto mayor sea la escala de generación, más fácil será que el modelo encuentre trayectorias incorrectas que «engañan al verificador»; si la false-negative rate
\[
\beta=P(v=0\mid z=1)
\]
es demasiado alta, se descartan soluciones correctas y diversas, y la distribución de entrenamiento se vuelve cada vez más estrecha. Aquí $\alpha$ controla la contaminación por errores y $\beta$ controla la cobertura de lo correcto; ambas deben reportarse por dificultad de la tarea y por slice de riesgo.

\lecturefigure{slide-25.jpg}{Un verifier confiable es la infraestructura central de RL reasoning}{25}

\begin{warningbox}{«La respuesta final es correcta» no garantiza que la trayectoria lo sea}
Un verifier de resultados puede aceptar trayectorias que acertaron por casualidad, que cometieron errores de unidades que luego se compensaron, que usaron información filtrada o que dependen de un estado de herramientas no reproducible. El entrenamiento de alto riesgo debe combinar un outcome verifier, process checks, detección de contaminación, un sandbox de ejecución y muestreo revisado por personas, en lugar de fijarse solo en la cadena final.
\end{warningbox}

\teachervoice{00:27:30--00:30:20, el docente enfatiza que, en las tareas que hoy pueden verificarse automáticamente, lo decisivo es el verifier, no la discusión sobre cuál actualización de RL es más elegante. Nota de clase: el límite superior de la escalabilidad lo determina «si es posible juzgar lo bueno y lo malo de forma barata y estable».}

\subsection{Qué escala realmente el entrenamiento de razonamiento}

La slide 26 vuelve a la teoría inicial: el razonamiento permite escalar la output length, mientras que la capacidad de responder directamente suele depender de la model depth. Al entrenar un reasoning model, la «scale» tiene al menos cuatro dimensiones: número de parámetros, número de problemas de entrenamiento, número de muestras por problema y longitud máxima de la trayectoria. Reportar solo los FLOPs totales oculta cómo se distribuyen estos presupuestos y tampoco permite explicar si la ganancia proviene de un modelo más fuerte, de más exploración o de un cómputo más largo.

\lecturefigure{slide-26.jpg}{Scaling reasoning: distinguir entre longitud de salida y profundidad del modelo}{26}

Si se muestrean $K$ trayectorias por problema con longitud promedio $T$, el costo de generación del entrenamiento es aproximadamente proporcional a $KT$; si la tasa de ejemplos aceptables es $q$, el costo esperado de obtener una trayectoria positiva es de alrededor de $KT/(Kq)=T/q$, pero aumentar $K$ eleva la probabilidad $1-(1-q)^K$ de «encontrar al menos una trayectoria correcta». Esto explica por qué las tareas con recompensa escasa consumen rápidamente el presupuesto de inferencia y muestra también que el verifier y la dificultad del currículo deben diseñarse en conjunto.

\begin{knowledgebox}{Términos clave: SFT, Rejection Sampling, RL Finetuning}
\textbf{SFT} imita secuencias de referencia dadas; \textbf{rejection sampling} muestrea del modelo actual y conserva solo las trayectorias que pasan la verificación; \textbf{RL finetuning}, de forma más general, optimiza directamente la recompensa esperada y puede aprovechar a la vez ejemplos positivos y negativos. Los tres pueden formar un pipeline continuo y no deben dividirse en bandos mutuamente excluyentes solo por el nombre de los artículos.
\end{knowledgebox}

\subsection{Resumen de la sección}

El giro central de RL finetuning es pasar de «reproducir una trayectoria escrita por una persona» a «aumentar, dentro de la distribución de candidatos del propio modelo, la probabilidad de las trayectorias exitosas según una calidad verificable». Esto proporciona una fuente cíclica de datos para escalar el razonamiento, pero también coloca en el centro del sistema los sesgos del verifier, el análisis de la respuesta, las fallas de la recompensa y el costo del muestreo.
